% Manual técnico de DevStudio
% Compilar: /usr/bin/latexbuild manual-tecnico.tex
% Preámbulo compartido por el reporte, los manuales y el contrato.
% Se compila con XeLaTeX mediante /usr/bin/latexbuild.
\documentclass[11pt,a4paper]{article}

\usepackage{fontspec}
\setmonofont{JetBrainsMono Nerd Font Mono}[Scale=0.82]
\usepackage[spanish,es-nodecimaldot,es-noshorthands]{babel}

\usepackage[top=2.4cm,bottom=2.4cm,left=2.4cm,right=2.4cm]{geometry}
\usepackage{fancyhdr}
\usepackage{titlesec}
\usepackage{xcolor}
\usepackage{booktabs}
\usepackage{tabularx}
\usepackage{array}
\usepackage{enumitem}
\usepackage{float}
\usepackage{graphicx}
\usepackage{tikz}
\usetikzlibrary{positioning,arrows.meta}
\usepackage{listings}
\usepackage{xurl}

\makeatletter
\renewcommand*\l@section[2]{%
  \ifnum \c@tocdepth >\z@
    \addpenalty\@secpenalty
    \addvspace{1.0em \@plus\p@}%
    \setlength\@tempdima{2.6em}%
    \begingroup
      \parindent \z@ \rightskip \@pnumwidth
      \parfillskip -\@pnumwidth
      \leavevmode \bfseries
      \advance\leftskip\@tempdima
      \hskip -\leftskip
      #1\nobreak\hfil \nobreak\hb@xt@\@pnumwidth{\hss #2}\par
    \endgroup
  \fi}
\makeatother
\usepackage[hidelinks]{hyperref}

\definecolor{tnBlue}{HTML}{0EA5E9}
\definecolor{tnDark}{HTML}{0F172A}
\definecolor{codeGray}{HTML}{F4F6F8}
\definecolor{codeGreen}{HTML}{2E7D32}
\definecolor{codeRed}{HTML}{C62828}

\newcolumntype{Y}{>{\raggedright\arraybackslash}X}

\lstset{
  basicstyle=\ttfamily\footnotesize,
  breaklines=true,
  breakatwhitespace=false,
  frame=tb,
  framerule=0.4pt,
  rulecolor=\color{tnDark!40},
  framesep=5pt,
  keywordstyle=\color{tnBlue}\bfseries,
  commentstyle=\color{codeGreen}\itshape,
  stringstyle=\color{codeRed},
  showstringspaces=false,
  tabsize=2,
  columns=fullflexible,
  keepspaces=true,
  extendedchars=true,
  captionpos=t,
  xleftmargin=0pt,
  xrightmargin=0pt,
  aboveskip=0.9em,
  belowskip=0.9em,
  literate=
    {á}{{\'a}}1 {é}{{\'e}}1 {í}{{\'i}}1 {ó}{{\'o}}1 {ú}{{\'u}}1
    {Á}{{\'A}}1 {É}{{\'E}}1 {Í}{{\'I}}1 {Ó}{{\'O}}1 {Ú}{{\'U}}1
    {ñ}{{\~n}}1 {Ñ}{{\~N}}1 {ü}{{\"u}}1 {Ü}{{\"U}}1
    {¿}{{?`}}1 {¡}{{!`}}1
    {·}{{\textperiodcentered}}1
}

\pagestyle{fancy}
\fancyhf{}
\fancyhead[L]{\small Alfredo Ramírez Candia}
\fancyhead[R]{\small Diseño de Aplicaciones Web}
\fancyfoot[C]{\thepage}
\renewcommand{\headrulewidth}{0.4pt}

\titleformat{\section}
  {\Large\bfseries\color{tnDark}}{\thesection.}{0.5em}{}
\titleformat{name=\section,numberless}
  {\Large\bfseries\color{tnDark}}{}{0em}{}
\titleformat{\subsection}
  {\large\bfseries\color{tnBlue}}{\thesubsection}{0.5em}{}
\titlespacing*{\section}{0pt}{1.4em}{0.6em}
\titlespacing*{\subsection}{0pt}{1.05em}{0.35em}

\addto\captionsspanish{%
  \renewcommand{\tablename}{Tabla}%
  \renewcommand{\figurename}{Figura}%
  \renewcommand{\lstlistingname}{Listado}%
}

\newcommand{\datosPdf}[1]{%
  \hypersetup{%
    pdftitle={#1},%
    pdfauthor={Alfredo Ramírez Candia},%
    pdflang={es}%
  }%
}

\newcommand{\portada}[3]{%
  \begin{titlepage}
    \centering
    \vspace*{0.8cm}
    {\large\bfseries BENEMÉRITA UNIVERSIDAD AUTÓNOMA DE PUEBLA\par}
    \vspace{0.28cm}
    {\large FACULTAD DE CIENCIAS DE LA COMPUTACIÓN\par}
    \vspace{0.28cm}
    {\large LICENCIATURA/INGENIERÍA EN CIENCIAS DE LA COMPUTACIÓN\par}
    \vspace{1.35cm}
    \rule{\linewidth}{0.5pt}\par
    \vspace{0.45cm}
    {\Huge\bfseries #1\par}
    \vspace{0.38cm}
    {\Large #2\par}
    \vspace{0.32cm}
    {\large Asignatura: Diseño de Aplicaciones Web\par}
    \vspace{0.45cm}
    \rule{\linewidth}{0.5pt}\par
    \vspace{1.5cm}
    \begin{tabular}{rl}
      \textbf{Docente:}   & Miguel Ángel Vargas Lomelí \\[0.55cm]
      \textbf{Alumno:}    & Alfredo Ramírez Candia \\[0.55cm]
      \textbf{Matrícula:} & 202375947 \\[0.55cm]
      \textbf{Fecha:}     & 22 de noviembre de 2026 \\
    \end{tabular}
    \vfill
    {\small #3\par}
  \end{titlepage}
  \setcounter{page}{2}
}

\setcounter{tocdepth}{2}
\setcounter{secnumdepth}{2}
\setlist[enumerate]{leftmargin=1.5em,itemsep=0.45em,topsep=0.4em}
\setlist[itemize]{leftmargin=1.4em,itemsep=0.28em,topsep=0.35em}
\emergencystretch=2em
\urlstyle{tt}
\renewcommand{\arraystretch}{1.18}

\datosPdf{Manual técnico. DevStudio}

\begin{document}
\portada{Manual técnico}{Instalación, API y base de DevStudio}{Manual técnico}

\tableofcontents
\clearpage

\section{Descripción}

DevStudio separa un frente estático de una API PHP. Nginx publica ambos en el
mismo origen. SQL Server guarda usuarios, catálogo, proyectos, avances,
reuniones, solicitudes y mensajes. Este manual permite instalar el sistema y
operar la API sin recorrer el resto de la documentación.

\section{Arquitectura}

\begin{figure}[H]
  \centering
  \begin{tikzpicture}[
      caja/.style={
        draw=tnDark!75, line width=0.7pt, rounded corners=3pt,
        fill=codeGray, align=center, inner sep=7pt, text width=8.6cm,
        font=\small
      },
      flecha/.style={-{Stealth[length=2.3mm]}, line width=0.75pt, draw=tnDark}
    ]
    \node[caja] (nav) {Navegador};
    \node[caja, below=7mm of nav] (ngx)
      {Nginx \texttt{pf-web}, puerto 8080};
    \node[caja, below=7mm of ngx] (php)
      {PHP 8.3 \texttt{pf-api}, red interna};
    \node[caja, below=7mm of php] (sql)
      {SQL Server \texttt{pf-mssql}, \texttt{DevStudioDB}};
    \draw[flecha] (nav) -- (ngx);
    \draw[flecha] (ngx) -- (php);
    \draw[flecha] (php) -- (sql);
  \end{tikzpicture}
  \caption{El navegador no abre un puerto distinto para la API.}
  \label{fig:arq}
\end{figure}

\path{appWeb/} se monta de solo lectura en Nginx. \path{backEnd/} se monta en
el contenedor PHP. Un cambio de HTML o CSS se ve al recargar. Cambiar
\path{nginx.conf} exige recrear el contenedor web. Recrear la API pierde las
sesiones PHP.

\section{Requisitos}

\begin{itemize}
  \item Docker Engine con Compose v2.
  \item Unos 2 GB libres para SQL Server.
  \item Arquitectura \texttt{x86\_64}. La imagen oficial de SQL Server no corre nativa en ARM.
  \item Puertos 8080 y 1433 libres, salvo que se cambien en \path{.env}.
  \item Cliente SQL opcional: \texttt{sqlcmd} ya viene dentro del contenedor.
\end{itemize}

\section{Instalación}

Desde \path{proyecto-final/}:

\begin{lstlisting}[caption={Arranque.}]
cp .env.example .env
docker compose up -d --build
\end{lstlisting}

\texttt{DB\_PASSWORD} es obligatorio. No se copia a los PDF ni al ZIP de
entrega. Después:

\begin{lstlisting}[caption={Esquema y datos iniciales.}]
docker compose exec -T mssql /bin/bash -c \
  '/opt/mssql-tools18/bin/sqlcmd -S localhost -U sa -P "$MSSQL_SA_PASSWORD" -C -i /dev/stdin' \
  < backEnd/database/schema.sql
docker compose exec -T mssql /bin/bash -c \
  '/opt/mssql-tools18/bin/sqlcmd -S localhost -U sa -P "$MSSQL_SA_PASSWORD" -C -i /dev/stdin' \
  < backEnd/database/seed.sql
\end{lstlisting}

Los dos scripts se pueden ejecutar más de una vez. \path{/api/v1/health} debe
responder en \url{http://localhost:8080/api/v1/health}.

Si la base ya existía de una etapa anterior, las migraciones
\path{001} a \path{005} en \path{backEnd/database/migrations/} agregan planes,
portafolio, perfil, progreso, reuniones y la columna de la foto sin borrar
datos.

\section{Configuración}

Las variables viven en \path{.env}. Las que cambian el comportamiento visible
son \texttt{WEB\_PORT}, \texttt{DB\_PORT}, \texttt{DB\_NAME},
\texttt{DB\_PASSWORD}, \texttt{APP\_URL} y los textos \texttt{COMPANY\_*}.
\texttt{SESSION\_NAME} identifica la cookie. En local,
\texttt{SESSION\_SECURE} queda en falso porque no hay HTTPS.

\section{Base de datos}

\begin{table}[H]
  \centering
  \caption{Tablas y la relación que importa.}
  \label{tab:tablas}
  \small
  \begin{tabularx}{\textwidth}{@{}l Y@{}}
    \toprule
    Tabla & Guarda \\
    \midrule
    \texttt{Roles} & Solo \texttt{admin} y \texttt{cliente} \\
    \texttt{Usuarios} & Cuenta, hash, rol y extensión de la foto \\
    \texttt{Clientes} & Empresa, contacto, cargo, frase y ciudad \\
    \texttt{Servicios}, \texttt{Tecnologias} & Catálogo público \\
    \texttt{Proyectos} & Trabajo, estado, publicación y \texttt{progreso} \\
    \texttt{ProyectoAvances} & Historial de porcentaje y nota \\
    \texttt{Reuniones} & Cita, estado, puntos y respuesta \\
    \texttt{Planes} & Recomendación del diagnóstico \\
    \texttt{SolicitudesCotizacion} & Pedidos de propuesta \\
    \texttt{MensajesContacto} & Bandeja de contacto \\
    \bottomrule
  \end{tabularx}
\end{table}

\texttt{Proyectos.progreso} no pasa de 100. Una reunión exige inicio menor que
el fin y una duración entre 15 y 240 minutos. El archivo de la foto no está en
la base: se escribe en \path{backEnd/storage/avatars/} y la columna
\texttt{foto} solo guarda \texttt{jpg}, \texttt{png} o \texttt{webp}.

\section{API}

La prefijo es \path{/api/v1}. Las respuestas de éxito van en \texttt{data}.
Un error trae \texttt{code}, \texttt{message} y, si aplica, \texttt{details}.

\subsection{Público y sesión}

\begin{table}[H]
  \centering
  \caption{Rutas sin rol de portal.}
  \small
  \begin{tabularx}{\textwidth}{@{}l Y@{}}
    \toprule
    Ruta & Uso \\
    \midrule
    \texttt{GET /health} & Estado del proceso y de la base \\
    \texttt{GET /config}, \texttt{/planes}, \texttt{/servicios}, \texttt{/portafolio} & Catálogo público \\
    \texttt{POST /contacto}, \texttt{POST /diagnostico} & Formularios, con límite por IP \\
    \texttt{GET /auth/csrf} & Token para la siguiente escritura \\
    \texttt{POST /auth/login}, \texttt{/registro}, \texttt{/logout} & Sesión \\
    \texttt{GET /auth/me} & Usuario actual \\
    \texttt{PATCH /auth/password} & Cambio de contraseña \\
    \bottomrule
  \end{tabularx}
\end{table}

El login devuelve un token CSRF nuevo. El token pedido antes de entrar no
sirve para las escrituras posteriores.

\subsection{Portal}

Todas exigen rol \texttt{cliente}. El identificador del cliente sale de la
sesión.

\begin{itemize}
  \item \texttt{GET /mi/resumen}, \texttt{/mi/proyectos} y \texttt{/mi/proyectos/\{id\}}.
        Los proyectos incluyen \texttt{progreso} y \texttt{avances}.
  \item \texttt{GET /mi/solicitudes}, detalle y \texttt{POST} para crear una.
  \item \texttt{GET} y \texttt{PATCH /mi/perfil}. El correo no se cambia aquí.
  \item \texttt{GET} y \texttt{POST /mi/foto}. El cuerpo es \texttt{multipart}
        con el campo \texttt{foto}.
  \item \texttt{GET /mi/reuniones?desde=\&hasta=}, \texttt{POST} para agendar
        y \texttt{PATCH /mi/reuniones/\{id\}} solo para cancelar.
\end{itemize}

\subsection{Administración}

Exigen rol \texttt{admin}. Hay CRUD de usuarios, clientes, servicios,
tecnologías y proyectos, y bandejas de solicitudes y mensajes. Dos rutas no
tienen página:

\begin{itemize}
  \item \texttt{PATCH /admin/proyectos/\{id\}/progreso} con \texttt{progreso}
        de 0 a 100 y \texttt{nota} de 2 a 500 caracteres.
  \item \texttt{GET /admin/reuniones} lista las citas.
        \texttt{PATCH} sobre cada identificador pasa de \texttt{solicitada}
        a \texttt{confirmada} o \texttt{rechazada}, y guarda
        \texttt{respuesta} y \texttt{puntos}.
        Una reunión \texttt{cancelada} responde 409.
\end{itemize}

Un cliente que llame estas rutas recibe \texttt{SIN\_PERMISOS} con HTTP 403,
antes incluso de revisar el CSRF.

\section{Seguridad}

\begin{itemize}
  \item Contraseñas con \texttt{password\_hash} de PHP. El seed no guarda la clave en claro.
  \item Consultas con parámetros. Los \texttt{\{id\}} de la ruta son dígitos.
  \item CSRF en las escrituras del portal, del panel y de la sesión.
  \item Cinco fallos de login en 15 minutos bloquean ese correo y esa IP.
  \item La foto se valida con \texttt{getimagesize}. No se acepta SVG.
  \item El directorio de avatares no lo publica Nginx.
\end{itemize}

\section{Frente}

El JavaScript del portal llama a la API con la cookie de sesión. La miniatura
de la foto se guarda en el navegador para pintarla antes de las hojas de
estilo. Cerrar sesión borra esa copia local. El calendario arma una cuadrícula
de lunes a domingo y traduce los instantes UTC a la hora de la Ciudad de México.

\section{Operación}

\begin{lstlisting}[caption={Parar y ver registros.}]
docker compose ps
docker compose logs api --tail 50
docker compose stop
\end{lstlisting}

Un respaldo de la base se toma con las herramientas de SQL Server sobre el
volumen \texttt{mssql\_data}. Borrar ese volumen obliga a cargar de nuevo
\path{schema.sql} y \path{seed.sql}. Las fotos subidas están en
\path{backEnd/storage/avatars/} y no forman parte del esquema.

\section{Fallos frecuentes}

\begin{itemize}
  \item \textbf{La API no arranca.} Falta \texttt{DB\_PASSWORD} o SQL Server todavía no pasó el healthcheck.
  \item \textbf{401 al escribir.} Hay que usar el \texttt{csrf\_token} que devolvió el login, no el anterior.
  \item \textbf{403 en una ruta \texttt{/admin/}.} La cuenta es de cliente. Es el resultado correcto.
  \item \textbf{422 al agendar.} Duración fuera de rango, inicio en el pasado o empalme con otra reunión activa.
  \item \textbf{La foto parpadea.} Recargar una vez después de subirla. Las páginas siguientes leen la miniatura local.
  \item \textbf{El HTML viejo sigue en pantalla.} Nginx marca las páginas como \texttt{no-store}. Una recarga forzada descarta la copia del navegador.
\end{itemize}

\section{Referencias internas}

El detalle largo de columnas y de cada campo de la API sigue en
\path{backEnd/docs/}. Este manual es el que debe bastar para instalar y para
probar las rutas nuevas de avance, reuniones y foto.

\end{document}
